\section*{Symbols at a glance}
\phantomsection
\addcontentsline{toc}{section}{Symbols at a glance}
\label{app:notation}

Each object is introduced where it first appears; this page is a place
to return to. Here $p_j$ is the $j$th prime. The lane of $p$ consists
of integers whose least prime factor is $p$.

\renewcommand{\arraystretch}{1.35}
\begin{longtable}{@{}p{.23\linewidth}p{.72\linewidth}@{}}
\toprule
Symbol & Meaning \\
\midrule
\endhead
$\Gold$ & The original Goldbach constant. \\
$C_g$ & A constant carrying a binary message $g$:
$g_j=0$ closes the lane of $p_j$; $g_j=1$ opens it. \\
$x_g(n)$ & The $n$th binary digit of $C_g$. \\
$\ell(n)$ & For $n\ge2$, the index of its least prime factor:
$\lpf(n)=p_{\ell(n)}$. \\
$\lambda(n)$ & The Liouville sign, $(-1)^{\Omega(n)}$;
$\Omega$ counts prime factors with multiplicity. \\
$\delta_j$ & The density of positions in the lane of $p_j$. \\
$S_g$ & The total density of disabled lanes. \\
$M_k(g)$ & The correlation of $k$ signs $1-2x_g(n)$, conditioned on pairwise coprime digit positions. \\
$K(k)$ & The limiting probability that $k$ sampled integers are pairwise coprime. \\
$\Phi_g,\Psi_g$ & The generating functions for conditioned and unconditioned correlations, respectively. \\
$W,\Comp$ & Independent prime and composite random factors;
$W\Comp$ has the exponential law of mean one. \\
$f,\SymCDF$ & The density of $W$ on $(0,\infty)$, extended to the complex plane,
and the symmetric cumulative distribution function derived from it. \\
$G_g$ & The cumulative distribution function from sampling equally spaced
digit positions with coprime start and step. \\
$G_N,Y_N$ & The clean and observed labelled coprimality tables in the finite recovery experiment. \\
$T_N$ & Effective signal $N(1-2\varepsilon_N)^2$, where $\varepsilon_N$
is the probability of flipping each pair answer. \\
$D(n)$ & The effective count of prime-power parts $p^{v_p(n)}$,
measured by the entropy of their logarithmic shares. \\
$B_k$ & The set of primes still present after deletion round $k$. \\
\bottomrule
\end{longtable}
\renewcommand{\arraystretch}{1}

\noindent All logarithms are natural unless a base is specified.
Local cutoff variables are defined where used.
